\begin{table*}[t]
\centering
\scriptsize
\caption{How this survey differs from adjacent survey traditions. The comparison concerns scope and evidence use, not the value of each literature.}
\label{tab:coverage-matrix}
\renewcommand{\arraystretch}{1.08}
\setlength{\tabcolsep}{2.4pt}
\begin{tabularx}{\textwidth}{@{}p{0.18\textwidth}Y Y Y Y@{}}
\toprule
\textbf{Dimension} & \textbf{Closest VLA safety survey} & \textbf{Embodied security survey} & \textbf{LLM-agent security survey} & \textbf{This survey} \\
\midrule
Direct VLA attack evidence & Discusses VLA threats within broader safety. & Usually sparse or absent. & Usually absent. & Centers demonstrated attacks on VLA policies, VLA-controlled robots, and VLA training/privacy channels. \\
Evidence use & Threat/evaluation oriented. & Vulnerability oriented. & Software-agent risk oriented. & Separates direct VLA evidence from embodied-adjacent and contextual work so adjacent risks do not inflate VLA claims. \\
Action representation & Often covered as model capability. & Usually not central. & Not applicable. & Compares action token, chunk, waypoint, trajectory, velocity/flow, skill/API, controller, and memory representations. \\
Controller/recovery & Defense or safety mechanism discussion. & Cyber-physical component discussion. & Not applicable. & Highlights why chunk horizon, controller frequency, safety filters, recovery, and pre/post-controller traces change attack claims. \\
Privacy attacks & Usually safety-adjacent. & Sometimes privacy/security. & Often data leakage. & Treats VLA membership inference and VLALeaks as direct VLA confidentiality evidence. \\
Evaluation emphasis & Broad trend synthesis. & Broad vulnerability synthesis. & Benchmark scores or attack rates. & Interprets typed attack success, validation tier, controller boundary, transfer, privacy metrics, and recovery rather than pooling incompatible ASR. \\
\bottomrule
\end{tabularx}
\end{table*}
